\chapter*{Executive Summary}
\addcontentsline{toc}{chapter}{Executive Summary}
\textbf{Problem and motivation.} Retail stores record CCTV footage mainly for security, yet managers lack a low-cost way to know whether customers are attended to, how staff time is distributed, and which zones receive traffic. Commercial retail-analytics products typically depend on dedicated sensors or hardware and on subscription costs, which puts them out of reach of many small and mid-sized retailers.

\textbf{Proposed solution.} RetailPulse reuses existing CCTV infrastructure. A coordinated set of AI agents (Detection, Classification, Engagement, Trajectory, Orchestrator, and Insight Narration), orchestrated with LangGraph, converts video into per-entity analytics: footfall, approach time, interaction duration, idle time, unattended-customer flags, dwell time, paths, and heatmaps. A POS analytics module adds product performance, purchase conversion, demand forecasts, and inventory recommendations. Results are presented on a web dashboard for managers and a focused workspace for employees.

\textbf{Scope.} The committed scope is a single-store prototype with three cameras, running on CPU only. Multi-store scaling, full cross-camera re-identification, shelf-level interaction detection, and a live POS connector are deliberately deferred to future work after a feasibility triage against hardware, timeline, and dependency constraints.

\textbf{Current status.} The Next.js frontend covers the manager, employee, and shared pages. The FastAPI and SQLite POS service supports CSV/XLSX import, deterministic metrics, 7- and 30-day forecasts, recommendations, and cited assistant answers. Role-aware access is enforced with signed HttpOnly session cookies and an internal secret header. An experimental camera-handoff pipeline links the cam1 right-side exit to the cam2 left-side entry. A local floor-distance CLI additionally reuses trained role detections and ByteTrack, stores editable provisional calibration, and exports supported floor positions/distances with motion-quality states. The vision agents and the LangGraph orchestration in the dashboard are still partly represented by prototype or demo data.

\textbf{Validation, risks, and remaining work.} Floor-distance verification includes 17 focused software checks, four existing handoff checks and actual sampled-video integration. Only the reference snapshot supported approximate floor projection; remaining sampled frames were motion-invalid or unverified. The main risks are CPU-only inference speed, role-classification accuracy, uncertain floor scale/foot proxies, camera movement and cross-camera association errors. The remaining work covers vision-agent integration, evaluation on store footage, orchestration, hardening, and the final demonstration.
